\section{Related Works}

\subsection{Articulated Object Reconstruction}

Articulated object reconstruction is a long-standing problem in 3D computer vision and has been widely researched in recent years. 
Feed-forward network-based methods~\cite{ditto,kawana2023detection,qian2022understanding,sun2024opdmulti,geng2023gapartnet,heppert2023carto,yu2024gamma,wang2024rpmart} are trained on an annotated dataset, but fail to generalize to unseen objects. 
To improve the generalization ability, a series of methods leverage the foundation models~\cite{articulateanything,realcode,chen2024urdformer,qiu2025articulate} to generate articulated object models from single-view images. While these approaches benefit from extensive pre-training on diverse datasets, they typically lack geometric consistency constraints and iterative refinement mechanisms. 
Another series regards articulated object reconstruction as calibrated multi-view camera capturing under two distinct articulation configurations~\cite{splart,reartgs,artgs,guo2025articulatedgs,liu2023paris,wang2024sm,weng2024neural}. Despite their geometric rigor, it remains difficult to align the axes of different states, leading to the limited practicality. 
To address these limitations and improve generalization, recently, a new line of methods~\cite{rsrd,video2articulation,wu2025predict,monomobility} reconstruct articulated objects from monocular RGB or RGB-D video sequences under the assumption that one part remains stationary relative to the background. 
% RSRD~\cite{rsrd} performs 4D reconstruction of articulated objects from monocular videos, though it necessitates an initial object reconstruction. Video2Articulation~\cite{video2articulation} builds interactable digital twins from casually captured RGB-D videos, yet similarly requires one part to be static relative to the background.
These methods suffer from two fundamental limitations. First, the static base part assumption is violated in practical scenarios where users naturally manipulate objects like scissors or pliers. Second, the inability to freely repose the object during capture results in incomplete coverage. In contrast, our method operates in a free-moving setting, where the object pose and joint state can vary concurrently, eliminating these issues.
\vspace{-2pt}
\subsection{Dynamic Reconstruction}
Articulated object reconstruction can be regarded as another type of dynamic reconstruction. Feed-forward methods~\cite{zhang2024monst3r,lu2025align3r,wang2025continuous} directly learn to reconstruct dynamic point clouds from large-scale datasets. However, they fail to recover precise motions under the free-moving setting (Fig.~\ref{fig:vis}). 
Optimization-based dynamic reconstruction methods~\cite{wu20244d,cao2023hexplane,huang2024sc,lin2024gaussian} reconstruct temporal deformations of radiance fields, but often lack generalization ability. Recently, point tracking methods~\cite{alltracker,tapip3d,xiao2025spatialtrackerv2,rajivc2025multi,ngodelta,yao2025uni4d} provide generalized priors at pixel-level resolution, enabling their application to free-moving articulated object reconstruction. However, due to the data-driven nature, the tracks inevitably contain noise and outliers. Our method uses an off-the-shelf point tracking model~\cite{alltracker} to generate pseudo motion labels, and optimizes the articulated object to fit the labels. In this way, we combine the generalization ability of point tracking and the high accuracy of the optimization-based dynamic reconstruction in one framework.